\section{Related Work}

\textbf{Multilingual and low-resource NER.} Cross-lingual transfer with
shared subword encoders substantially narrowed the gap between
high-resource and low-resource named-entity recognition
\citep{oyelaran2022lowresource}, but transfer quality still tracks the
morphological distance between source and target languages
\citep{ferreira2019attention}. Span-based formulations of joint entity and
relation extraction \citep{tanaka2020span} generalize more gracefully to
new entity types than sequence-tagging formulations, which is part of why
we adopt a span detector rather than a BIO tagger in \sys{VertexTag}.

\textbf{Entity normalization and linking.} Normalization is often treated
as a downstream step of entity linking, where a candidate generator
proposes canonical forms and a reranker scores them against context
\citep{schneider2020adapter}. Most prior candidate generators are built for
Wikipedia-scale knowledge bases; the AfricaNLP shared task instead supplies
smaller, curated gazetteers, which changes the trade-off between recall
(favoring aggressive fuzzy matching) and precision (favoring a strong
reranker) that a system must strike.

\textbf{Discourse and document-level consistency.} Because the same entity
often appears multiple times in a document under different surface forms,
normalization quality is tied to document-level coherence modeling
\citep{nakamura2021survey}. We do not model discourse explicitly in
\sys{VertexTag}, but our error analysis in Section~\ref{sec:analysis}
suggests that document-level re-scoring of low-confidence normalizations is
a promising direction for future shared-task submissions.

\textbf{Evaluation under drift.} Automatic metrics for text generation and
extraction tasks are known to be sensitive to distribution shift between
development and test splits \citep{kapoor2022evaluation,
ibrahim2023calibration}. The AfricaNLP organizers mitigated this by drawing
the test split from a later time window than training and development data;
Section~\ref{sec:task-data} describes the resulting split in detail.
